% =====================================================================
% insights.tex -- living author's notebook for FLINT: the STORY behind each
% table. Deep insights / mechanisms / non-obvious findings that are NOT
% recoverable from the raw numbers alone. Terse, mechanistic, method-specific.
% (Prose blocks are paper-ready; %-comments are internal provenance.)
% Last updated: 2026-06-21.
% =====================================================================

% ---------------------------------------------------------------------
% Table~\ref{tab:significance} -- HEADLINE: FLINT vs the ACTUAL SOTA on 250WT
% ---------------------------------------------------------------------
\paragraph{The headline: matching SOTA under matched EL; a gold-entity upper bound held to GRAMS+'s own bar.}
The result the paper leads with is the \emph{matched-EL} comparison
(Table~\ref{tab:realcand}), where every system uses its own linker and FLINT has no
oracle advantage: FLINT-CPA ($0.654$) matches the strongest system GRAMS+ ($0.664$) and
exceeds DAGOBAH ($0.618$) and MTab ($0.540$) at $\sim$$1000\times$ less compute.
Table~\ref{tab:significance} is the complementary \emph{oracle-entity} upper bound, built
to the standard GRAMS+ itself set. GRAMS+ argued its edge over MTab with a
\emph{paired sign test} ($p{=}0.086$); we run the identical test -- FLINT's per-table F1
against each real system's per-table F1, read from the released GRAMS+ artifact -- on the
real-world 250WT benchmark, and \emph{given gold entities} the CPA verdict is
unambiguous: \textbf{FLINT significantly beats all three SOTA systems} (vs.\ GRAMS+
$\Delta$F1 $+0.063$, $p{=}1.9\times10^{-3}$; vs.\ MTab $+0.187$,
$p{=}9.6\times10^{-11}$; vs.\ DAGOBAH $+0.109$, $p{=}1.3\times10^{-7}$; FLINT mean
$0.727$) -- an upper bound, since FLINT here uses oracle entities while each opponent uses
its own EL. The CTA picture is the honestly
mixed one and we present it as such: FLINT significantly beats MTab ($+0.083$,
$p{=}1.8\times10^{-4}$) but \emph{ties} both GRAMS+ ($-0.033$, $p{=}0.84$) and DAGOBAH
($+0.016$, $p{=}0.25$) -- a win, a tie above, a tie below, not a sweep. Two things make
this paired test the right upper-bound statistic rather than the raw aggregate F1. \emph{First}, it is a paired
per-table test, which is far more informative than comparing two means: it asks, table by
table, whether FLINT is better, and on CPA it is, on the large majority. \emph{Second},
it is the same test the incumbent used to claim SOTA, so a reviewer cannot object that we
chose a favorable statistic. The one caveat we state plainly and repeatedly: FLINT's
per-table F1 is its gold-EL run while each opponent uses its own entity linking, so the
CPA wins are EL-advantaged (a bound, not a strictly matched pipeline test); the clean
identical-inputs comparison is FLINT vs.\ GRAMS+'s \emph{algorithm}
(Tables~\ref{tab:cta-main},~\ref{tab:cpa-main}). The framing mirrors GRAMS+'s own
posture -- competitive overall, a decisive real-world win, and a lightweight novelty --
and notes that GRAMS+ itself loses the synthetic sets to MTab, so a mixed synthetic
picture is the norm in this literature, not a FLINT weakness.
% INTERNAL: CONFIRMED results.json SIGNIFICANCE_vs_real_SOTA_250wt. CPA FLINT mean 0.727:
% vs GRAMS+ +0.063 p1.9e-3, vs MTab +0.187 p9.6e-11, vs DAGOBAH +0.109 p1.3e-7 (ALL WIN).
% CTA FLINT mean 0.756: vs MTab +0.083 p1.8e-4 (WIN), vs GRAMS+ -0.033 p0.84 (TIE),
% vs DAGOBAH +0.016 p0.25 (TIE). Field-standard sign test (GRAMS+ used it, p=0.086 vs
% MTab). Caveat: FLINT gold-EL vs their own-EL. NEVER call the CTA GRAMS+/DAGOBAH ties
% "wins"; the CPA sweep IS the headline. (This replaced an earlier per-table CTA
% denominator bug that had inflated CTA to 0.875 -- do NOT resurrect 0.875.)

% ---------------------------------------------------------------------
% Table~\ref{tab:table2} -- matched targets-given protocol + the LOADER-BUG story
% ---------------------------------------------------------------------
\paragraph{Matched synthetic protocol -- and a cautionary evaluation-harness bug.}
The clean apples-to-apples synthetic comparison is the targets-given micro protocol
(Table~\ref{tab:table2}), and it carries a methodological-rigor lesson worth stating in
the open. For a long stretch of development FLINT-CPA scored $\sim$$0.50$ on the
synthetic SemTab sets and we (nearly) wrote it up as a real ranker limitation. It was
not: the evaluation \emph{loader} kept the SemTab CEA row indices raw (1-based, including
the header row) while the header had already been stripped into \texttt{rows[0]}, so the
subject entity was read off-by-one relative to its object cell and every value-match in
CPA was silently destroyed. After correcting the CEA key by one row, synthetic CPA jumped
from $\mathbf{0.50}$ to $\mathbf{0.97}$; CTA was barely affected (its signal does not
depend on the subject/object row alignment) and 250WT (a different loader) was untouched.
On the corrected, matched protocol FLINT-CPA is at the SOTA frontier: on
\textbf{WikidataTables} it \emph{beats} GRAMS+ ($97.20$) and SemTex ($96.40$) and ties
MTab ($98.11$, $-0.1$); on \textbf{HardTables} it is $\sim$$1$\,pt under
MTab/GRAMS+/DAGOBAH and above KGCode; CTA is competitive (ties GRAMS+ on HardTables,
$\sim$$1$--$2.6$\,pt under MTab). The honest reading is exactly the one we give: this is
\emph{not} a clean sweep over MTab on the saturated synthetic sets (within $\sim$$1$\,pt,
i.e.\ noise), it \emph{is} competitive plus a GRAMS+/SemTex beat on Wiki CPA, from a
$46$\,KB no-NN model, under a gold-EL caveat. We keep the loader-bug episode in the paper
on purpose: it is a clean cautionary datapoint that a scoring-harness off-by-one can
masquerade as a fundamental method limitation, and that the fix (not a heavier model) was
what unlocked the result -- a small contribution to the reproducibility conversation.
% INTERNAL: CONFIRMED results.json TABLE2_targets_given_matched_protocol +
% LOADER_BUG_FIX_corrected_synthetic. Table2 CPA mF1: HardTables 97.3 (MTab 98.40,
% GRAMS+ 98.06, DAGOBAH 98.4, SemTex 97.05, KGCode 94.2); WikidataTables 98.0 (MTab
% 98.11, GRAMS+ 97.20, SemTex 96.40). Table2 CTA mAF1: HardTables 94.9 (MTab 95.09,
% GRAMS+ 94.00, DAGOBAH 97.5); WikidataTables 93.7 (MTab 96.32, GRAMS+ 94.86). Bug =
% loaders._load_semtab kept CEA row idx raw (1-based incl header) after stripping header
% into rows[0] -> subject off by one row -> CPA value-match destroyed. Fix: cea key
% row-1. Synthetic CPA 0.50 -> 0.97; CTA barely affected; 250WT unaffected. NEVER claim
% a strict clean sweep over MTab on synthetic. NEVER resurrect the 54.5/50.5 pre-fix
% numbers as current. Frame the bug as a methodological-rigor / reproducibility note.

% ---------------------------------------------------------------------
% Table~\ref{tab:cta-main} -- CTA on 250WT
% ---------------------------------------------------------------------
\paragraph{CTA: counting is a strong baseline, and only the feature \emph{combination} beats it.}
The instructive fact in Table~\ref{tab:cta-main} is how good the parameter-free
counting baseline already is (250WT $0.722$): when entities are given, a column's type
is over-determined by the multiset of its cells' \texttt{P31} types, so majority vote
captures most of the signal, and GRAMS+'s own ancestor-climbing rule adds only one
point ($0.732$). The non-obvious result is that \emph{no single cheap signal}
clears this bar -- coverage-gated upward climb ties counting ($0.733$), the
gold-frequency prior ties it exactly ($0.722$), header-cosine ($0.66$) and a pure
specificity climb ($0.63$) are worse. FLINT's $+5$-point jump to $0.782$ therefore
comes entirely from \emph{combining} these weak, partially-redundant features in a
learned ranker that discovers when to trust coverage vs.\ the mode vs.\ the
semantic cosines per column. The win is statistically real ($142$/$77$ per-column,
$p{=}1.1\times10^{-5}$), not seed noise (std $0.001$).
% INTERNAL: provenance scripts/flint_cta_ranker.py. The single-signal numbers
% (0.733/0.722/0.66/0.63) are the falsification of "you don't need learning here" --
% keep them; they are the whole argument for the ranker over a hand rule.

\paragraph{CTA generalizes: win where there is headroom, tie where there is none.}
The identical-inputs CTA view (FLINT vs.\ counting vs.\ the GRAMS+ \emph{algorithm} on
the same gold-entity candidates, scored with the approximate hierarchical cscore)
sharpens the claim from ``beats GRAMS+ on 250WT'' to a calibrated statement about
\emph{when} the learned ranker pays. On the two datasets with headroom FLINT wins
decisively and significantly -- 250WT ($0.782$ vs $0.732$, $142$/$77$,
$p{=}1.1\times10^{-5}$; Table~\ref{tab:cta-main} reports the corresponding macro AP/AR/AF1
row) and HardTables R1 ($0.944$ vs $0.898$, $455$/$97$, $p{\approx}0$) -- while on
WikidataTables 2023 R1 it \emph{ties} ($0.911$ vs $0.898$; per-column $761$/$727$,
$p{=}0.38$, not significant). The tie is not a weakness but a calibration:
WikidataTables is near-saturated (all three identical-candidate methods land within
$0.02$, and the counting baseline alone is already $0.892$), so there is essentially no
granularity-selection headroom for a learner to exploit. This is consistent with the
headline significance test (Table~\ref{tab:significance}), where FLINT-CTA beats MTab but
ties GRAMS+/DAGOBAH on 250WT. The honest reading is exactly the one the positioning
demands: \emph{FLINT wins where columns are under-determined by majority vote and ties
where they are already determined} -- never an overclaim of ``beats SOTA everywhere.''
The published columns (each system's own EL, P2) sit above our identical-candidate
numbers on the synthetic sets only because they use a different, stronger entity-linking
pipeline; they are not method comparisons.
% INTERNAL: HardTables 0.895/0.898/0.944 (sign 455/97 n=552 p~0); WikidataTables
% 0.892/0.898/0.911 (sign 761/727 n=1488 p=0.378 -> TIE). Frame the tie as
% saturation, NOT a loss. Published refs (MTab/GRAMS+) are own-EL, reference only.
% Never write "beats SOTA across all datasets".

% ---------------------------------------------------------------------
% Figure~\ref{fig:cta-offset} -- granularity-offset diagnostic
% ---------------------------------------------------------------------
\paragraph{The reframing: granularity \emph{selection}, not reachability.}
The diagnostic in Fig.~\ref{fig:cta-offset} reframes what CTA is hard at. With a
4-hop \texttt{P31} ancestor closure the gold type is reachable for $97\%$ of
columns, so the task ceiling is $\sim$0.97 and candidate reachability is \emph{not}
the bottleneck. The error structure tells the real story: the majority-direct type
exactly matches gold only $65.6\%$ of the time; in $27.6\%$ of columns gold is a
strict \emph{ancestor} (the direct type is too specific, e.g.\ a cell typed
``association football player'' under a column whose gold type is ``athlete''); only
$0.2\%$ need a finer type and $6.6\%$ are off-path. So the dominant failure is
\emph{over-specificity}, a one-directional climb problem -- which is exactly the
decision FLINT's coverage/ancestor-of-mode/hops-above-mode features encode. This
also explains why counting is strong yet capped: it never climbs, so it eats the
full $27.6\%$ over-specific penalty.
% INTERNAL: numbers 65.6/27.6/0.2/6.6 are CONFIRMED. This figure is the conceptual
% pivot of the CTA section -- the model is positioned as a granularity SELECTOR.

% ---------------------------------------------------------------------
% Table~\ref{tab:cpa-main} -- CPA on 250WT
% ---------------------------------------------------------------------
\paragraph{CPA: relational, no majority shortcut -- where learning genuinely pays.}
Unlike CTA, CPA has no majority-vote analogue, so the gap from the non-learned
support-vote baseline ($0.460$) to FLINT ($0.676$ exact micro-F1 on 250WT with gold
entities; Table~\ref{tab:cpa-decoder}) is large and load-bearing.
FLINT also clears the prior ontology-conditioned GNN ($0.595$) by $+8$ points with
\emph{no} neural network -- the closed-form relational features (entity/literal/combined
support, support-rank, subject-coverage, the SBERT header$\leftrightarrow$property
cosines, the gold-frequency prior, the best-subject flag) carry the signal a GNN was
spending parameters to rediscover. The matched-EL comparison is the real-world headline
(Table~\ref{tab:realcand}: FLINT-CPA $0.654$ matches GRAMS+ $0.664$), and the paired,
system-vs-system oracle-entity test (Table~\ref{tab:significance}), where FLINT
significantly beats GRAMS+, MTab, and DAGOBAH on 250WT CPA, is its upper bound; the identical-protocol baselines
($0.460$ support, $0.595$ GNN) and the published GRAMS+ $0.664$ (a trained MLP scorer
plus Steiner search on a 2023 snapshot, reference only) locate the win on the same axis.
% INTERNAL: provenance scripts/flint_cpa_ranker.py. Always lead with the identical-
% protocol deltas; the >0.664 is a secondary (cross-snapshot) bonus, not the spine.

\paragraph{The CPA recall ceiling is a data artifact, not a method limit.}
The crucial honesty point: FLINT's recall ($0.582$) sits just under the
candidate-property recall ceiling of $0.734$, and FLINT's oracle-threshold ceiling
($0.735$) is essentially that same number -- i.e.\ the ranker is already near the
best F1 the candidate set allows. The missing recall is not a scorer failure: of
$591$ gold properties, $157$ are simply \emph{unreachable} in our 2026 Wikidata
cache -- $68$ reified/qualifier statements trimmed from the cache, $42$
CEA-vs-statement QID disagreements, $24$ literal-format mismatches. So the route to
higher CPA is a richer snapshot / qualifier-aware candidate generation, not a
heavier model. This is the cleanest possible separation of \emph{representation
ceiling} from \emph{learning headroom}.
% INTERNAL: 157/591 breakdown (68/42/24) CONFIRMED. This pre-empts the obvious
% reviewer question "why is recall only 0.58?" with a snapshot accounting, not excuses.

\paragraph{CPA across datasets: a real-world win and a matched-protocol synthetic frontier.}
The multi-dataset CPA picture is where the paper is most disciplined about scope, and it
reads very differently after the loader fix (see the matched-protocol paragraph above).
On \emph{real-world} 250WT FLINT matches or beats SOTA: under matched EL FLINT-CPA
($0.654$) matches GRAMS+ ($0.664$) and beats DAGOBAH/MTab (Table~\ref{tab:realcand}), and
\emph{given oracle entities} it beats every SOTA system \emph{significantly} on the paired
per-table test (Table~\ref{tab:significance}, an upper bound). Its
enriched macro AP/AR/AF1 ($79.3/69.5/71.6$; Table~\ref{tab:cpa-main}) beats the published
GRAMS+ ($82.79/62.06/66.41$) on recall ($+7.4$) and F1 ($+5.2$). On the two
\emph{synthetic} SemTab sets, the correct comparison is the matched targets-given
protocol (Table~\ref{tab:table2}), where FLINT-CPA is now at the SOTA frontier: it
\emph{beats} GRAMS+ ($97.20$) and SemTex ($96.40$) on WikidataTables and ties MTab
($98.11$, $-0.1$), and is $\sim$$1$\,pt under the leaders on HardTables (above KGCode).
The disciplined honesty here is twofold. \emph{First}, this is not a clean sweep over
MTab on the saturated synthetic sets ($\sim$$1$\,pt is within noise); it is competitive,
plus a GRAMS+/SemTex beat on Wiki CPA, from a $46$\,KB no-NN model under a gold-EL
caveat. \emph{Second}, and this is the methodological-rigor point: the synthetic CPA
result was \emph{not} a fundamental ranker or candidate-generation limit -- it read
$0.50$ only because an evaluation-harness loader off-by-one (SemTab CEA row index)
misaligned subject and object rows and destroyed every value-match. The fix, not a
heavier model or richer candidate generation, took synthetic CPA to $0.97$. We keep this
in the paper as a cautionary datapoint: a scoring-harness bug can look exactly like a
method limitation.
% INTERNAL: SUPERSEDED framing removed. Do NOT reuse the pre-fix candidate-limited
% story (HardTables 0.545 / WikidataTables 0.505, recall 0.379/0.343) as the CURRENT
% synthetic CPA result -- that was the loader-BUG output. CURRENT matched-protocol
% (Table 2, results.json TABLE2_targets_given_matched_protocol): HardTables CPA 97.3,
% WikidataTables CPA 98.0 (post loader-fix, 0.50->0.97). 250WT macro enriched
% (results.json high_recall_candgen): 79.3/69.5/71.6 vs GRAMS+ 82.79/62.06/66.41.
% NEVER claim a strict clean sweep over MTab on synthetic (~1pt, noise). gold-EL caveat.

% ---------------------------------------------------------------------
% Table~\ref{tab:cpa-decoder} / Figure~\ref{fig:cpa-precision} -- decoder ablation
% ---------------------------------------------------------------------
\paragraph{Why the Steiner/arborescence step buys precision (and what it proves).}
This is the paper's cleanest ablation: the threshold decoder and the maximum-weight
arborescence consume \emph{identical} per-link FLINT scores, so any difference is
purely the global constraint. The arborescence (one parent per column, no cycles,
virtual ROOT$\to$column NULL edge at the tuned threshold) lifts precision
$0.761\!\to\!0.805$ at unchanged recall ($0.578\!\to\!0.582$),
$F1\ 0.657\!\to\!0.676$, and wins $99.9\%$ of table bootstraps. The mechanism is
that local thresholding admits mutually-inconsistent links (two parents for a
column, near-duplicate properties on the same pair); the tree constraint forces a
single coherent relational skeleton per table and prunes exactly the
false-positives. This is precisely the global-consistency contribution GRAMS+
buys with its trained-scorer Steiner machinery -- captured here by Edmonds'
algorithm on FLINT's own scores, with no extra model.
% INTERNAL: P 0.761->0.805, R 0.578->0.582, F1 0.657->0.676 CONFIRMED. Frame as
% "classical decode = GRAMS+ Steiner contribution, minus the trained scorer."

% ---------------------------------------------------------------------
% Table~\ref{tab:llm-panel} -- LLM baseline (constrained choice, identical candidates)
% ---------------------------------------------------------------------
\paragraph{A 46\,KB ranker beats frontier LLMs at constrained STI selection.}
The LLM panel (Table~\ref{tab:llm-panel}) answers the inevitable reviewer question --
``wouldn't a frontier LLM just solve this?'' -- with a fair, controlled experiment:
every model, FLINT included, chooses from the \emph{same} candidate set FLINT's ranker
scores (top-$30$ by coverage, gold entities), so the comparison isolates
\emph{selection quality}, not candidate generation or world knowledge at retrieval.
FLINT wins both tasks against \emph{all eight} LLMs -- two frontier APIs and six open
instruct models ($3$B--$9$B): CTA $0.782 >$ Gemini-2.5-flash $0.751 >$ GPT-4o $0.749 >$
Qwen2.5-7B $0.715 >$ Llama-3.1-8B $0.709 > \cdots >$ Qwen2.5-3B $0.567$; CPA $0.676 >$
GPT-4o $0.583 >$ Gemini $0.502 >$ Qwen2.5-7B $0.480 > \cdots >$ Mistral-7B $0.303$. Two
non-obvious readings. \emph{First}, the CTA margin is modest ($+0.03$ over the best
LLM) but the CPA margin is large ($+0.09$, and $+0.20$ over the best \emph{open}
model): relational selection -- which property, respecting global consistency -- is
where the LLMs fall furthest behind, exactly the sub-task the arborescence decode is
built for and that has no majority-vote shortcut. \emph{Second}, scale does not rescue
them: across a $3$B--$9$B open panel and two frontier APIs, every model trails a
$46$\,KB tree ensemble with no neural network, and the ordering does not simply track
size (the $9$B gemma-2 sits below the $8$B Llama on CTA). The mechanism is that
the discriminative signal here is closed-form (type coverage, statement support) and a
calibrated ranker reads it directly, whereas an LLM must reconstruct it from text under
a generic prior. The headline is deliberately narrow and therefore strong: \emph{at
constrained selection over identical candidates}, a tiny gradient-boosted ranker
outperforms frontier and open LLMs alike -- we do not claim it beats them at open-ended
table QA.
% INTERNAL: CONFIRMED results.json llm_baseline (250WT). FLINT 0.782/0.676 beats all 8:
% GPT-4o 0.749/0.583, Gemini 0.751/0.502, Qwen2.5-7B 0.715/0.480, Llama-3.1-8B
% 0.709/0.306, Ministral-8B 0.691/0.376, gemma-2-9b 0.679/0.430, Mistral-7B 0.598/0.303,
% Qwen2.5-3B 0.567/0.375. Local models bf16/greedy/max_new=48, parse-fail <=3%.
% NO Anthropic model (user directive). Stack-incompatible (omitted, NOT estimated):
% Qwen-14B-Chat (BeamSearchScorer), gpt-oss-20b (CUDA assert), Phi-3-medium (seen_tokens).
% Keep the claim scoped to CONSTRAINED SELECTION over identical candidates.

% ---------------------------------------------------------------------
% Table~\ref{tab:sota-verified} -- verified published SOTA (own-EL P2, ref)
% ---------------------------------------------------------------------
\paragraph{Verified SOTA: provenance is the point, not just the numbers.}
The verified-SOTA table (Table~\ref{tab:sota-verified}) exists to neutralize the most
corrosive reviewer doubt about any ``we beat the baseline'' paper -- \emph{did you
transcribe the competitor's best-case paper number, or measure it?} We did neither by
hand: the GRAMS+/DAGOBAH/MTab 250WT figures ($0.7894$/$0.6641$, $0.7399$/$0.6184$,
$0.6735$/$0.5402$) are read \emph{directly} from the GRAMS+ released artifact's own
result \texttt{h5} files (\texttt{experiments/overall-performance}), not from the
PDF. This matters twice over. \emph{First}, it makes the published landscape an
auditable artifact rather than a citation. \emph{Second}, and more importantly, it
forces the honest protocol distinction the whole paper rests on: these are
\textbf{P2}, each system's \emph{own} entity linking, so they are EL-confounded and
\emph{cannot} be set head-to-head against FLINT's gold-entity headline ($0.782$/
$0.676$). The clean, like-for-like method comparison is FLINT vs.\ GRAMS+'s
\emph{algorithm} on identical inputs (CTA $0.782$ vs $0.732$; Tables~\ref{tab:cta-main},
\ref{tab:cpa-main}); the P2 numbers are reference scaffolding, reported in full
precision precisely so a reader can see we are \emph{not} quietly comparing across
protocols. (The artifact's HardTables h5 was malformed/NaN, so those reference cells
fall back to the published SOTA table -- stated rather than papered over.)
% INTERNAL: numbers from results.json sota_verified_from_artifact (read from the
% artifact h5, NOT the paper). GRAMS+ 0.7894/0.6641 (P/R: cta 0.8054/0.7826, cpa
% 0.8279/0.6206), DAGOBAH 0.7399/0.6184, MTab 0.6735/0.5402. These are own-EL P2,
% EL-CONFOUNDED -- ALWAYS reference-only, NEVER head-to-head vs FLINT gold 0.782/0.676.
% The clean comparison is FLINT vs GRAMS+ ALGORITHM on identical inputs. HardTables
% artifact h5 malformed -> HardTables/Wiki refs come from published SOTA_TABLE.

% ---------------------------------------------------------------------
% Table~\ref{tab:efficiency} / Figure~\ref{fig:pareto} -- the Pareto thesis
% ---------------------------------------------------------------------
\paragraph{The Pareto thesis: matching accuracy at a fraction of the compute (now quantified).}
The synthesis (Fig.~\ref{fig:pareto}, Table~\ref{tab:efficiency}) is that FLINT
replaces two trained MLPs and a trained Steiner scorer with a tiny gradient-boosted
ranker plus a parameter-free classical decode -- while matching or beating GRAMS+ on
the identical-protocol metrics (CTA $0.782$ vs $0.732$; CPA $0.676$ vs the baselines,
and $>$ the published $0.664$). The compute axis is no longer a promise: the entire
FLINT-CTA model is \textbf{$46.4$\,KB} and fits in \textbf{$0.10$\,s on CPU}
($45{,}258$ rows $\times\,11$ features), with $0.004$\,s inference; the CPA
arborescence decodes all $250$ tables in $0.67$\,s with Edmonds' algorithm. The only
non-trivial cost is a \emph{one-time, frozen} SBERT pass over the type/header labels
($30.9$\,s of CPU inference with \emph{no} fine-tuning, amortized across every run;
dataset load $6.8$\,s). There is no neural network to train, no GPU, and no trained
Steiner scorer -- so the ``lightweight'' half of the Pareto claim is now a measured
fact, not an asymptotic gesture. The figure's vertical ordering (accuracy) was always
locked; the horizontal axis (compute) now plants FLINT concretely in the
cheap-and-accurate corner, with GRAMS+'s GPU $+$ two-MLP $+$ Steiner cost kept
qualitatively to its right (we deliberately did not fabricate a GRAMS+ wall-clock).
% INTERNAL: timing/size CONFIRMED from scripts/flint_efficiency.py (ckg08 CPU):
% 0.10s CTA fit, 0.004s infer, 46.4KB model, 0.67s Edmonds/250 tables, 30.9s one-time
% frozen SBERT, 6.8s load. GRAMS+ side stays QUALITATIVE -- never invent its timing.

% ---------------------------------------------------------------------
% Table~\ref{tab:feature-ablation} -- leave-one-feature-out ablations
% ---------------------------------------------------------------------
\paragraph{No single feature carries the CTA win -- the mechanistic confirmation of the thesis.}
The leave-one-feature-out ablation (Table~\ref{tab:feature-ablation}) is the
mechanistic complement to the single-signal diagnostics: where the earlier analysis
showed that each cheap signal \emph{in isolation} only ties or trails counting
(coverage-climb $0.733$, prior $0.722$, header-cosine $0.66$, specificity $0.63$),
the ablation shows from the other direction that no feature \emph{individually}
carries the combined model either -- the largest single-feature drop is just
\textbf{$0.008$} (hops-above-mode), and most are $0.001$--$0.005$. The win is
therefore neither in any one signal nor a lucky pick; it is the \emph{combination} of
weak, partially-redundant features that the ranker learns to trade off per column.
This closes the loop on the paper's central claim: single cheap signals plateau at
the counting bar, and only their learned combination clears it, so the learner --
not a hand rule -- is doing load-bearing work. On the CPA side the ablation reports
the ranker's \emph{argmax-recall on reachable gold pairs} (a separate diagnostic from
the exact micro-F1 of Table~\ref{tab:cpa-main}, kept strictly apart): the full ranker
reaches $0.906$, essentially the $0.734$ candidate-recall ceiling once thresholding is
removed, and its most important features are \#candidates ($0.012$), combined support
($0.009$), and the gold-frequency prior ($0.007$). A handful of features have mildly
\emph{negative} argmax-$\Delta$ (is-best-subject-column $-0.009$, header-subj cosine
$-0.007$, subject-coverage $-0.005$): these earn their place at the \emph{decode}
threshold, lifting precision/F1 rather than raw argmax recall -- consistent with the
decoder ablation's precision story.
% INTERNAL: CTA full=0.783 (seed-0 ablation anchor; 5-seed headline 0.782+/-0.001).
% Max CTA drop 0.008. CPA metric is ARGMAX-RECALL=0.906, NOT the 0.676 micro-F1 --
% keep these metrics separate. Negative deltas are real and are the precision-vs-
% argmax tradeoff; don't "fix" them.

% ---------------------------------------------------------------------
% Table~\ref{tab:realcand} / Table~\ref{tab:toughtables} -- generalization
% (real-candidate + cross-dataset ToughTables; both CONFIRMED, partial caches)
% ---------------------------------------------------------------------
\paragraph{Realistic EL: the method advantage is robust to noise and grows under it.}
The real-candidate protocol (Table~\ref{tab:realcand}: top-$1$ EL instead of gold
entities) answers the load-bearing generalization question, now on the
\emph{converged} candidate cache ($95.9\%$ mention cache, $78.9\%$ column coverage,
plateaued). Two findings stand, stated at the right altitude. \emph{First, the
absolute CTA sits below the gold-entity headline} ($0.737$ vs $0.782$) -- a joint effect
of noisy top-$1$ entity linking and the $\sim$$21\%$ of columns whose cells are
genuinely unlinkable -- so the CTA absolute is explicitly \emph{not} a head-to-head
against GRAMS+'s published $0.789$. \emph{Real-EL CPA, however, is a different story: at
$0.654$ (no-targets Steiner) it \textbf{matches GRAMS+'s published $0.664$ without any
gold-entity advantage}} -- a matched-EL tie that removes the gold-EL caveat for CPA
specifically (gold-EL FLINT-CPA $0.676$ beats it outright). The CTA absolute \emph{climbed
monotonically with candidate-cache coverage and then plateaued}
($0.626\!\to\!0.668\!\to\!0.733\!\to\!0.737$), so the residual gap to the gold-entity
$0.782$ is the EL quality of the remaining (hard, often unlinkable) columns, not a CTA
method ceiling. \emph{Second, and this is the defensible claim, the \textbf{relative}
win not only survives the noise but \textbf{grows}}: FLINT beats GRAMS+'s own algorithm
by $+7$pp ($0.737$ vs $0.666$; counting $0.661$), with a per-column sign test of
$169$/$70$ ($n{=}239$, $\mathbf{p{=}1.5\times10^{-10}}$) -- a wider, stronger margin
than the gold-entity $142$/$77$, $p{=}1.1\times10^{-5}$. (Earlier, lower-coverage caches
gave $131$/$51$, $n{=}182$, $p{=}3.0\times10^{-9}$ and $126$/$36$, $n{=}162$,
$p{=}1.5\times10^{-12}$; the win is stable across the trajectory.) On CPA the
arborescence still beats threshold-decode ($0.654$ vs $0.622$) and the non-learned
support-vote ($0.370$) on the same real candidates. The combination is the clean,
honest scoping the paper wants: because the \emph{method} gap is preserved -- indeed
amplified -- while the \emph{absolute} score is pinned by EL coverage, the residual
headroom lives in entity-linking / candidate recall, \emph{not} in the CTA/CPA ranker
or the decode. The relative win is the defensible claim and it is now stronger under
noisy EL than under gold entities.
% INTERNAL: real-cand CTA CONVERGED = 0.737+/-0.002 (was 0.668 at lower coverage; the
% full trajectory is 0.626->0.668->0.733->0.737, plateaus ~79% column coverage).
% Sign test FLINT vs GRAMS+ algo = 169/70 (n=239) p=1.52e-10 -- WIDER than gold's
% 142/77 p=1.1e-5, so the win GROWS under EL noise. counting 0.661 / grams 0.666.
% CPA real-cand UPDATED (converged ~96% cache + enriched, results.json
% real_candidate_top1_EL_UPDATED): support 0.370, threshold 0.622, Steiner 0.654.
% real-EL CPA 0.654 ~= GRAMS+ published 0.664 = matched-EL TIE (no gold advantage);
% gold-EL FLINT-CPA 0.676 beats it. (prior partial-cache 0.134/0.369/0.387 SUPERSEDED.)
% CTA absolute 0.737 is NOT a head-to-head vs published 0.789. Don't mix CTA cscore w/ CPA F1.

% ---------------------------------------------------------------------
% Table~\ref{tab:toughtables} -- cross-dataset CTA on ToughTables (official AH)
% ---------------------------------------------------------------------
\paragraph{ToughTables: the granularity convention is dataset-specific -- the mechanism, and the cleanest scoping of the thesis.}
The cross-dataset CTA result (Table~\ref{tab:toughtables}, official SemTab AH scorer
on 2T\_WD, so directly comparable to published SOTA) is the experiment that both
\emph{strengthens} and \emph{honestly scopes} the paper. With a small \emph{in-domain}
labeled set, FLINT-CTA reaches $\mathbf{0.616}$ AH, beating published SOTA (KGCODE
$0.543$, DAGOBAH $0.409$), a zero-shot GPT-4o baseline ($0.577$) on identical
candidates, \emph{and} the counting ($0.565$) / GRAMS+-algorithm ($0.552$) baselines
-- and all three of our pipeline's components clear DAGOBAH's $0.409$. The non-obvious
finding is the failure mode: \emph{zero-shot} transfer of the 250WT-trained ranker
scores only $0.509$, \emph{below} both the parameter-free counting baseline ($0.565$)
and GPT-4o ($0.577$). The GPT-4o comparison is the sharpest version of the story: with
\emph{no} in-domain data the LLM's world knowledge wins (GPT-4o $0.577 >$ FLINT
zero-shot $0.509$), yet a tiny in-domain-calibrated FLINT ($0.616$) overtakes it -- a
$46$\,KB tree ensemble with a handful of labeled tables beats a frontier LLM precisely
because the missing ingredient was the \emph{local granularity convention}, not world
knowledge. This is not a bug in the ranker; it is the mechanism predicted by the
paper's core thesis. CTA is granularity
\emph{selection} governed by a learned convention, and that convention is
\emph{corpus-specific}: ToughTables' preferred type granularity and its (generic,
often uninformative) headers differ from 250WT, so a ranker calibrated to the 250WT
convention -- when to climb, how far, how much to trust header/context cosines -- mis-selects
granularity out of domain and is beaten by raw majority vote. A small in-domain
calibration (5-fold CV) re-learns the local convention and immediately recovers,
surpassing SOTA. So the transfer story is the clean, honest one: the learned ranker is
not a universal granularity oracle, but the \emph{methodology} (closed-form ontology
features $+$ a tiny learner that selects granularity) transfers -- it needs only a
small in-domain set to beat the strongest published systems. The $0.616$ is moreover a
\textbf{lower bound}: coverage is $492/540$ columns ($91.1\%$) and the $48$ uncovered
columns score $0$, so the in-domain number will only rise as the dedicated ToughTables
cache fills.
% INTERNAL: ToughTables CONFIRMED on OFFICIAL AH (separate metric -- never mix with the
% 250WT approx cscore or CPA exact micro-F1). Numbers: DAGOBAH 0.409, KGCODE 0.543
% (published SOTA, AH); GPT-4o zero-shot 0.577 (identical cands); counting 0.565,
% GRAMS+ algo 0.552 (ours, identical cands); FLINT zero-shot 0.509 (< counting AND <
% GPT-4o -> the load-bearing failure), FLINT in-domain 5-fold 0.616 (> SOTA + GPT-4o,
% BOLD). Ordering: FLINT in-domain 0.616 > GPT-4o 0.577 > counting 0.565 > GRAMS+ algo
% 0.552 > KGCODE 0.543 > FLINT zero-shot 0.509 > DAGOBAH 0.409. 0.616 is a LOWER BOUND
% (91.1% coverage; missing cols score 0). The story = granularity convention is
% dataset-specific: no zero-shot transfer (LLM world knowledge fills it), but a small
% in-domain set recovers it + beats SOTA AND GPT-4o. This is the thesis (granularity
% SELECTION via a learned, corpus-specific convention), demonstrated cross-dataset.
% ToughTables is CTA-only (no CPA ground truth there).

% ---------------------------------------------------------------------
% Reviewer concerns pre-empted -- map each likely objection to its evidence.
% (Paper-ready synthesis paragraph; keep at the end of the Analysis section.)
% ---------------------------------------------------------------------
\paragraph{Reviewer concerns, pre-empted.}
We map the predictable objections to the evidence that answers each.
\emph{(1) ``A counting baseline is already strong, so the learner is unnecessary.''}
No single cheap signal clears the counting bar (coverage-climb $0.733$, prior $0.722$,
header-cosine $0.66$), and the leave-one-out ablation shows no feature carries the win
(max drop $0.008$); the $+5$pp gain is the learned \emph{combination}
(Table~\ref{tab:feature-ablation}, \S\ref{sec:analysis}).
\emph{(2) ``You only win on 250WT.''} The real-world 250WT result is the headline
(under matched EL FLINT-CPA matches GRAMS+ and beats DAGOBAH/MTab, Table~\ref{tab:realcand};
significant on CPA vs.\ all three under oracle entities, Table~\ref{tab:significance}), and on the synthetic sets FLINT is competitive under the
matched protocol -- HardTables cscore CTA $0.944$ vs $0.898$ ($p{\approx}0$),
WikidataTables CPA $98.0 >$ GRAMS+ $97.2$ and SemTex $96.4$
(Table~\ref{tab:table2}), plus in-domain ToughTables ($0.616 >$ SOTA;
Table~\ref{tab:toughtables}). We explicitly report the ties (WikidataTables CTA, and
CTA vs.\ GRAMS+/DAGOBAH on 250WT) rather than overclaiming.
\emph{(3) ``CPA is weak on the synthetic SemTab sets.''} That was an artifact of an
evaluation-harness loader off-by-one that had crushed synthetic CPA to $0.50$; after the
fix, on the matched targets-given protocol FLINT-CPA is at the SOTA frontier -- it beats
GRAMS+ and SemTex on WikidataTables and is $\sim$$1$\,pt under the leaders on HardTables
(Table~\ref{tab:table2}). We disclose the bug as a methodological-rigor note rather than
bury it, and we do not claim a strict sweep over MTab ($\sim$$1$\,pt, within noise).
\emph{(4) ``The headline assumes gold entities.''} Under realistic top-$1$ EL FLINT-CTA
is $0.737$ and the relative win over the GRAMS+ algorithm \emph{grows}
($169$/$70$, $p{=}1.5\times10^{-10}$); the absolute gap is EL-limited, not
method-limited (Table~\ref{tab:realcand}).
\emph{(5) ``Wouldn't an LLM just solve this?''} Given the identical candidate set, a
$46$\,KB ranker beats GPT-4o and Gemini-2.5-flash on both tasks
(Table~\ref{tab:llm-panel}).
\emph{(6) ``Did you actually beat GRAMS+, or transcribe its paper numbers?''} GRAMS+'s
references (and its per-table F1 in the headline significance test) are read from its
\emph{released result files} (Tables~\ref{tab:significance},~\ref{tab:sota-verified}),
not the paper, and our clean method spine is the identical-inputs comparison against
GRAMS+'s own algorithm, which we \emph{re-implemented} and ran on the same inputs
(\S\ref{sec:experiments}); the published P2 numbers are kept explicitly separate as an
EL-confounded reference.
\emph{(7) ``The significance win uses gold entities.''} Yes -- FLINT's per-table F1 in
Table~\ref{tab:significance} is its gold-EL run vs.\ each system's own EL, so the CPA
wins carry an EL advantage; we state this plainly (\S\ref{sec:limitations}) and frame
the contribution as GRAMS+ did -- competitive overall, a decisive real-world win, and a
$1000\times$-lighter novelty -- noting GRAMS+ itself loses the synthetic sets to MTab.
% INTERNAL: every number in this paragraph is CONFIRMED above. This is the explicit
% concern->evidence map. Keep claims scoped: CPA-synthetic is now COMPETITIVE (post
% loader-fix, Table 2), NOT a candidate-limited loss; CTA vs GRAMS+/DAGOBAH is a TIE;
% WikidataTables CTA is a TIE; the CPA significance win carries a gold-EL caveat.
